% section: 漸化式を解く前に

\section{漸化式を解く前に}\label{節：漸化式を解く前に}

\subsection{数列と漸化式}\label{節：数列と漸化式}
数を順に並べたものを\bookterm{sequence}{数列}という.各項を$a_1,a_2,a_3,\ldots$と書き,添字で位置を表す.

例えば,$a_1=1$から始めて,直前の項に$2$を足すと,
\[
  1,\ 3,\ 5,\ 7,\ \ldots
\]
という\bookterm{sequence}{数列}ができる.この規則を式で書けば,
\[
  a_{n+1}=a_n+2\qquad(n\ge1)
\]
となる.項どうしの関係によって\bookterm{sequence}{数列}を定める式を,\bookterm{recurrence}{漸化式}という.
最初に与える値は\bookterm{initial}{初期条件}と呼ばれる.

次の項が一意に計算できる限り,同じ\bookterm{initial}{初期条件}と規則から同じ\bookterm{sequence}{数列}ができる.
直前の二項を使う規則では,通常,初めの二項を与える.
分母が零になるなど,途中で計算できなくなる場合は,その範囲を確かめる.

番号$n$から$a_n$を直接求める式が\bookterm{general}{一般項}である.
上の\bookterm{sequence}{数列}なら$a_n=2n-1$となる.
では,直前の項から次の項を求める式を使って,番号$n$から直接$a_n$を求めるにはどうすればよいだろうか.
まずは簡単な漸化式で,実際に計算してみよう.


\subsection{更新を追う}

  元の項をそのまま追うより,別の量を追うと次への変化が簡単になることがある.
  例えば$a_{n+1}=2a_n+3$では,定数$3$があるため等比型の式をそのまま使えない.
  両辺に$3$を加えると
  \[
    a_{n+1}+3=2(a_n+3)
  \]
  となる.$b_n=a_n+3$を追えば,更新は2倍するだけになる.目標は次のような形である.
  \begin{itemize}
    \item $b_{n+1}-b_n=f(n)$なら,変化を足し戻す（\bookterm{difference}{階差}型）.
    \item $b_{n+1}=f(n)b_n$なら,倍率を掛け合わせる（\bookterm{ratio}{階比}型）.
    \item $b_{n+1}=b_n$なら,その量は初項から変わらない.
    \item 既知の公式と同じ形なら,その公式の入力に当たる量を追う.
  \end{itemize}
  本書では,更新を追いやすい形を「きれいな形」と呼ぶ.和や積の表示をどこまで計算できるかは,別に確認する.
  基本4型の後に,\sectionref*{節：追う量を変える}と\sectionref*{節：解を組み合わせる見方}を学ぶ.
  元の式を観察し,追う量を作り,\bookterm{general}{一般項}や\bookterm{periodic}{周期},\bookterm{limit}{極限}を調べる.
  結論を元へ戻し,初期条件と更新式に合うことを確かめよう.
  \BookFigApproach
